% !TeX root = ../main.tex

% The word "Abstract" should be centered 2 inches below the top of 
% the page. Then, skip one line, then center your name followed by 
% the title of your thesis/dissertation. Use as many lines as necessary.
% Centered below the title include the phrase, in parentheses,
% "(Under the direction of _________)" and include the name(s) of 
% the members of your dissertation committee.
%
% Skip one line, then begin the content of the abstract. It should be 
% double-spaced and conform to margin guidelines. An abstract should not 
% exceed 150 words for a thesis and 350 words for a dissertation. The 
% latter is a requirement of both the Graduate School and UMI's 
% Dissertation Abstracts International.
%
% Because your dissertation abstract will be published, please prepare and 
% proofread it carefully. Print all symbols and foreign words clearly and 
% accurately to avoid errors or delays. Make sure that the title given at 
% the top of the abstract has the same wording as the title shown on your 
% title page. Avoid mathematical formulas, diagrams, and other 
% illustrative materials, and only offer the briefest possible description 
% of your thesis/dissertation and a concise summary of its conclusions. Do 
% not include lengthy explanations and opinions.
%
% The abstract should bear the lower case Roman number ii (if you did not 
% include a copyright page) or iii (if you include a copyright page).
\pretocheader{\bfseries\MakeTextUppercase{Abstract}}
\vspace{-1.5\baselineskip}

% Header must include your name in all caps, then the title of your
% dissertation/thesis in title font (not in all caps). Both of these
% elements must be center-aligned and single-spaced.
\begin{singlespace}
\begin{center}
\docAuthorFull: \docTitle\\
(Under the direction of Dr. Harvey Seim)
\end{center}
\end{singlespace}
\vspace{-0.5\baselineskip}

% The rest of the text can be double-spaced and left-aligned.
Ocean dynamics along coastlines are vastly influenced by topographic features. Particularly, The interactions between shoals and its surrounding flow play an important role in understanding the mixing in the global ocean. Despite this, there remains uncertainty about the role these shoals play in transporting fluid offshore and how they spur vertical and horizontal mixing. In this work, we use a series of large eddy simulations to simulate flow over shoals, analyzing the offshore transport of momentum and tracers, and the generation of eddy kinetic energy upstream and downstream. Our simulations are centered around a reference case of the Diamond Shoals off of North Carolina, using boundary conditions prescribed from observations. Our results show that the shoal serves as an extension of the coastline, diverting northward flow offshore and generating turbulence that enhances mixing. By testing variations in shoal geometry, we find that offshore transport and mixing are driven by the vertical and horizontal extent of the shoal bathymetry relative to the surrounding shelf. Additionally, shoals alter the water column depth and produce a hydraulic jump, transitioning supercritical upstream flow to subcritcal downstream flow. 
\clearpage
